\documentclass[11pt,a4paper]{article}
\usepackage[utf8]{inputenc}
\usepackage[T1]{fontenc}
\usepackage{amsfonts}
\usepackage[french]{babel}
\frenchbsetup{StandardLists=true}
\usepackage{enumitem}
\usepackage{amssymb}
\usepackage{amsmath}
\usepackage[left=2cm,right=2cm,top=2cm,bottom=2cm]{geometry}
\usepackage{multirow}
\usepackage[dvipsnames]{xcolor}

\usepackage[T1]{fontenc}
\usepackage{fouriernc}
\usepackage{graphicx}

\usepackage{setspace}
\singlespacing

\usepackage{pstricks}
\usepackage{qrcode}
\usepackage{xkeyval}
\usepackage{pst-labo}

\usepackage{tcolorbox}
\usepackage{multicol}
\usepackage{caption}
\usepackage{fancyhdr}
\pagestyle{fancy}
\renewcommand\headrulewidth{1pt}
\fancyhead[L]{Lycée Denis Diderot }
\fancyhead[R]{Classe de 3PM}
\setcounter{page}{1}\fancyfoot[C]{page \thepage}
\fancyfoot[R]{Année 2025-2026}
\fancyfoot[L]{Alexis RUIZ}
\usepackage{multirow,tabularx,booktabs}
\newcommand{\cadreblanc}[1]{%
	\medskip\framebox[\linewidth][c]{%
		\rule{0mm}{#1mm}}\par
	\medskip}

\usepackage{awesomebox}
\setlength{\aweboxvskip}{0mm}
\usepackage{pifont}
\usepackage{chemmacros}
\chemsetup{modules=all}
\setlength{\columnseprule}{.25mm}
\setlength{\parindent}{0pt}

\renewcommand{\arraystretch}{2}
\newcommand*\acocher{\quitvmode{\fboxsep0pt \fboxrule0.8pt \fbox{\vrule height1.8ex depth.1ex width0pt \vrule height0pt depth0pt width1.9ex }}}

\frenchbsetup{StandardLists=true}

\begin{document}
	
	\begin{center}
		\LARGE{Chapitre 4 - Les ions - EXERCICES}
		
		
		
	\end{center}
	
	\normalsize
	
	\hrule\
	\\[0,2cm]
	
	\textbf{Exercice 1 : Vrai ou Faux} \\
	
	Indiquer si les affirmations suivantes sont vraies ou fausses. 
	
	\begin{enumerate}
		\item Un ion est un atome qui a gagné ou perdu des protons. \acocher
		
		
		\item Un cation est un ion chargé positivement. \acocher
		
		
		\item L'ion chlorure \ch{Cl-} a gagné un électron. \acocher
		
		
		
		\item Une solution ionique peut être chargée électriquement. \acocher
		
		
		\item L'atome de sodium possède 11 protons et l'ion sodium \ch{Na+} possède 12 protons. \acocher
		
		
		
	\end{enumerate}
	
	
	
	\textbf{Exercice 2 : Atome ou ion ?} \\
	
	Compléter le tableau suivant :
	
	\vspace{5mm}
	
	\begin{center}
		\begin{tabular}{|l|c|c|c|c|}
			\hline
			\textbf{Particule} & \textbf{Nombre de protons} & \textbf{Nombre d'électrons} & \textbf{Atome ou ion ?} & \textbf{Formule} \\
			\hline
			Sodium & 11 & 11 & & \\
			\hline
			& 11 & 10 & & \\
			\hline
			Chlore & 17 & 17 & & \\
			\hline
			& 17 & 18 & & \\
			\hline
			& 29 & 27 & & \ch{Cu^{2+}} \\
			\hline
			Fer & 26 & 26 & & \\
			\hline
			& 26 & 23 & & \\
			\hline
		\end{tabular}
	\end{center}
	
	\vspace{5mm}
	
	
	
	\textbf{Exercice 3 : Cation ou anion ?} \\
	
	Pour chaque ion, indiquer s'il s'agit d'un cation ou d'un anion, puis préciser combien d'électrons ont été perdus ou gagnés par rapport à l'atome.
	
	\begin{enumerate}
		\item Ion potassium \ch{K+}
		
		\cadreblanc{15}
		
		\item Ion zinc \ch{Zn^{2+}}
		
		\cadreblanc{15}
		
		\item Ion sulfate \ch{SO4^{2-}}
		
		\cadreblanc{15}
		
		\item Ion fer (III) \ch{Fe^{3+}}
		
		\cadreblanc{15}
		
		\item Ion hydroxyde \ch{HO-}
		
		\cadreblanc{15}
	\end{enumerate}
	
	\vfill
	
	\textbf{Exercice 4 : Composition des ions} \\
	
	L'atome de fer possède 26 protons et 26 électrons.
	
	\begin{enumerate}
		\item Combien de protons et d'électrons possède l'ion fer (II) \ch{Fe^{2+}} ?
		
		\cadreblanc{15}
		
		\item Combien de protons et d'électrons possède l'ion fer (III) \ch{Fe^{3+}} ?
		
		\cadreblanc{15}
		
		\item Expliquer pourquoi ces deux ions ne portent pas la même charge.
		
		\cadreblanc{15}
	\end{enumerate}
	
	\vfill
	
	\newpage
	
	\textbf{Exercice 5 : Tests de reconnaissance} \\
	
	On réalise des tests sur trois solutions notées A, B et C. Voici les résultats obtenus :
	
	\begin{itemize}
		\item Solution A + hydroxyde de sodium → précipité bleu
		\item Solution B + hydroxyde de sodium → précipité vert
		\item Solution C + nitrate d'argent → précipité blanc qui noircit à la lumière
	\end{itemize}
	
	\vspace{5mm}
	
	\begin{enumerate}
		\item Quel ion est présent dans la solution A ? Justifier.
		
		\cadreblanc{20}
		
		\item Quel ion est présent dans la solution B ? Justifier.
		
		\cadreblanc{20}
		
		\item Quel ion est présent dans la solution C ? Justifier.
		
		\cadreblanc{20}
	\end{enumerate}
	
	\vfill
	
	\textbf{Exercice 6 : Identifier le bon test} \\
	
	On dispose de 4 solutions contenant respectivement des ions cuivre (II), fer (II), fer (III) et zinc. Pour chaque solution, indiquer :
	
	\begin{enumerate}[label=\alph*)]
		\item Le réactif à utiliser pour identifier l'ion
		\item La couleur du précipité obtenu
	\end{enumerate}
	
	\vspace{5mm}
	
	\textbf{Solution 1 :} Ions cuivre (II) \ch{Cu^{2+}}
	
	\cadreblanc{20}
	
	\textbf{Solution 2 :} Ions fer (II) \ch{Fe^{2+}}
	
	\cadreblanc{20}
	
	\textbf{Solution 3 :} Ions fer (III) \ch{Fe^{3+}}
	
	\cadreblanc{20}
	
	\textbf{Solution 4 :} Ions zinc \ch{Zn^{2+}}
	
	\cadreblanc{20}
	
	\vfill
	
	
	
	\textbf{Exercice 7 : Électroneutralité des solutions} \\
	
	Une solution ionique est toujours électriquement neutre.
	
	\vspace{5mm}
	
	\textbf{1.} Solution de chlorure de sodium \ch{NaCl}
	
	\begin{enumerate}[label=\alph*)]
		\item Quels sont les ions présents dans cette solution ?
		
		\cadreblanc{15}
		
		\item Quelle est la charge de chaque ion ?
		
		\cadreblanc{15}
		
		\item Vérifier que la solution est électriquement neutre.
		
		\cadreblanc{15}
	\end{enumerate}
	
	\vspace{5mm}
	
	\textbf{2.} Solution de sulfate de cuivre (II) \ch{CuSO4}
	
	\begin{enumerate}[label=\alph*)]
		\item Quels sont les ions présents dans cette solution ?
		
		\cadreblanc{15}
		
		\item Quelle est la charge de chaque ion ?
		
		\cadreblanc{15}
		
		\item Vérifier que la solution est électriquement neutre.
		
		\cadreblanc{15}
	\end{enumerate}
	
	\vfill
	
	\newpage
	
	\textbf{Exercice 8 : Formules des composés ioniques} \\
	
	Trouver la formule chimique des composés ioniques formés par les ions suivants :
	
	\vspace{5mm}
	
	\textbf{1.} Ions sodium \ch{Na+} et ions chlorure \ch{Cl-}
	
	\cadreblanc{15}
	
	\textbf{2.} Ions cuivre (II) \ch{Cu^{2+}} et ions chlorure \ch{Cl-}
	
	\cadreblanc{15}
	
	\textbf{3.} Ions sodium \ch{Na+} et ions sulfate \ch{SO4^{2-}}
	
	\cadreblanc{15}
	
	\textbf{4.} Ions zinc \ch{Zn^{2+}} et ions hydroxyde \ch{HO-}
	
	\cadreblanc{15}
	
	\textbf{5.} Ions fer (III) \ch{Fe^{3+}} et ions chlorure \ch{Cl-}
	
	\cadreblanc{15}
	
	\vspace{5mm}
	
	\begin{tcolorbox}[colback=red!5!white,colframe=red!75!black]
		{
			\textbf{Méthode :} Pour trouver la formule d'un composé ionique, il faut que la charge totale positive soit égale à la charge totale négative.
			
			\textbf{Exemple :} Ions aluminium \ch{Al^{3+}} et ions chlorure \ch{Cl-}
			
			Pour équilibrer les charges, il faut 3 ions \ch{Cl-} (3 fois -1 = -3) pour 1 ion \ch{Al^{3+}} (+3).
			
			La formule est donc : \ch{AlCl3}
		}
	\end{tcolorbox}
	
	\vfill
	
	\newpage
	
	\textbf{Exercice 9 : Problème - L'eau de mer} \\
	
	L'eau de mer contient principalement les ions suivants (concentrations en g/L) :
	
	\begin{itemize}
		\item Ions sodium \ch{Na+} : 10,8 g/L
		\item Ions chlorure \ch{Cl-} : 19,4 g/L
		\item Ions magnésium \ch{Mg^{2+}} : 1,3 g/L
		\item Ions sulfate \ch{SO4^{2-}} : 2,7 g/L
	\end{itemize}
	
	\vspace{5mm}
	
	\begin{enumerate}
		\item Parmi ces ions, quels sont les cations ? Quels sont les anions ?
		
		\cadreblanc{20}
		
		\item L'eau de mer a-t-elle un goût salé ? Justifier en citant le composé ionique responsable de ce goût.
		
		\cadreblanc{20}
		
		\item Quel test pourrait-on réaliser pour vérifier la présence d'ions chlorure dans l'eau de mer ? Décrire le résultat attendu.
		
		\cadreblanc{25}
		
		\item Expliquer pourquoi l'eau de mer conduit le courant électrique alors que l'eau pure ne le conduit pas.
		
		\cadreblanc{25}
	\end{enumerate}
	
	\vfill
	
	\newpage
	
	\textbf{Exercice 10 : Situation concrète - Les piles} \\
	
	Les piles électriques fonctionnent grâce au déplacement d'ions. Dans une pile alcaline (pile classique AA ou AAA), on trouve notamment des ions zinc \ch{Zn^{2+}} et des ions hydroxyde \ch{HO-}.
	
	\vspace{5mm}
	
	\begin{enumerate}
		\item L'ion zinc est-il un cation ou un anion ? Justifier.
		
		\cadreblanc{20}
		
		\item L'atome de zinc possède 30 protons. Combien d'électrons possède l'ion zinc \ch{Zn^{2+}} ?
		
		\cadreblanc{20}
		
		\item L'ion hydroxyde \ch{HO-} est-il un cation ou un anion ?
		
		\cadreblanc{15}
		
		\item Les ions hydroxyde et les ions zinc forment ensemble un composé ionique. Quelle est la formule de ce composé ? (Aide : penser à l'électroneutralité)
		
		\cadreblanc{25}
	\end{enumerate}
	
	\vfill
	
	\textbf{Exercice 11 : Recherche documentaire} \\
	
	À l'aide de tes connaissances ou en effectuant une recherche, réponds aux questions suivantes :
	
	\vspace{5mm}
	
	\begin{enumerate}
		\item Citer deux boissons du commerce qui contiennent des ions. Quels ions y trouve-t-on principalement ?
		
		\cadreblanc{25}
		
		\item Pourquoi ajoute-t-on du sel sur les routes en hiver ? Quel lien avec les ions ?
		
		\cadreblanc{25}
		
		\item Qu'est-ce qu'une eau « dure » ? Quels ions sont responsables de la dureté de l'eau ?
		
		\cadreblanc{25}
	\end{enumerate}
	
	\vfill
	
\end{document}